\ticket{24}{Машинный перевод: основные технологии автоматического перевода текста, их ограничения. Статистические и нейросетевые переводчики, принципы их построения}

\begin{examplan}
    \item Перечислить основные технологии машинного перевода.
    \item Описать правиловый, статистический и нейросетевой перевод.
    \item Раскрыть принципы построения SMT и NMT.
    \item Назвать ограничения технологий машинного перевода.
\end{examplan}

\subsection*{Короткая версия для письменного ответа}

\textbf{Машинный перевод} --- автоматическое преобразование текста с одного языка на другой с сохранением смысла. Основные подходы: перевод на правилах, статистический перевод, нейросетевой перевод и гибридные системы.

\subsection*{Перевод на правилах}

\textbf{RBMT} использует лингвистические правила, словари и грамматики.

Варианты:
\begin{itemize}
    \item прямой перевод: слово в слово с локальными перестановками;
    \item трансферный перевод: анализ исходного текста, преобразование структуры, генерация текста на целевом языке;
    \item интерлингва: перевод через универсальное семантическое представление.
\end{itemize}

Ограничения: огромная трудоемкость правил, хрупкость, слабая адаптация к новым доменам, трудности с идиомами и неформальным языком.

\subsection*{Статистический машинный перевод}

\textbf{SMT} обучается на параллельных корпусах и ищет наиболее вероятный перевод.

По Байесу:
\[
e^*=\arg\max_e P(e\mid f)=\arg\max_e P(f\mid e)P(e),
\]
где $P(f\mid e)$ --- модель перевода, а $P(e)$ --- языковая модель целевого языка.

Основные варианты:
\begin{itemize}
    \item word-based SMT;
    \item phrase-based SMT;
    \item syntax-based SMT.
\end{itemize}

Phrase-based SMT строится через выравнивание слов, извлечение пар фраз, оценку вероятностей фраз, модель переупорядочивания, языковую модель и декодер. Декодирование часто использует лучевой поиск.

Ограничения SMT: зависимость от больших параллельных корпусов, слабое моделирование дальних зависимостей, сложность настройки множества компонентов.

\subsection*{Нейросетевой машинный перевод}

\textbf{NMT} напрямую моделирует $P(e\mid f)$ нейронной сетью, обычно в архитектуре encoder-decoder.

Развитие архитектур:
\begin{itemize}
    \item RNN/LSTM/GRU encoder-decoder;
    \item attention, позволяющий декодеру обращаться к разным частям исходного предложения;
    \item Transformer, современный стандарт на self-attention.
\end{itemize}

Построение NMT:
\begin{enumerate}
    \item собрать параллельный корпус;
    \item выполнить токенизацию и разбиение на подслова, например BPE или SentencePiece;
    \item выбрать архитектуру и гиперпараметры;
    \item обучить модель end-to-end по кросс-энтропии;
    \item при переводе использовать авторегрессионную генерацию и beam search;
    \item оценивать качество метриками BLEU, METEOR, TER и ручной оценкой.
\end{enumerate}

Ограничения NMT: требовательность к данным и ресурсам, трудная интерпретация, ошибки на редких словах, падение качества на длинных предложениях, возможные галлюцинации и проблемы с терминологией.

\subsection*{Гибридные системы}

Гибридные подходы комбинируют правила, статистические модели и нейросети: например, используют правила для предобработки терминологии или постредактирования результата NMT.

\begin{important}
    \item RBMT основан на правилах и словарях.
    \item SMT выбирает перевод по вероятностной модели перевода и языковой модели.
    \item NMT обучает одну нейросетевую модель перевода end-to-end.
    \item Transformer-based NMT --- современный стандарт.
    \item Все подходы ограничены данными, неоднозначностью, контекстом, терминологией и оценкой качества.
\end{important}
